Independent increments and Gaussian structure reduce many Brownian calculations to covariance integrals. Deterministic square-integrable kernels give Gaussian It\^o integrals; exponentiation or evaluation at an independent random time generally does not preserve normality. Use orthogonal projection for jointly Gaussian observations and reflection for path maxima. Conditioning on an endpoint, a probability-zero event, requires densities or a continuous version of the conditional law.

An independent exponential failure time becomes a discount factor; a generator boundary problem can then be verified with It\^o's formula. Interchanging integrals, expectations or parameter derivatives requires nonnegativity, integrability or a suitable domination argument. Poisson splitting requires independent marking. In the renewal model, uniform observation over a long time interval asymptotically selects intervals in proportion to their lengths.
\par\smallskip
